\subsection{权与根格}

令 \(l = \dim V\)。记 \(Q(R)\) 为 \(V\) 中由 \(R\) 生成的子群；\(Q(R)\) 的元素称为 \(R\) 的根格元素。由第 6 号的 Th. 3，\(Q(R)\) 是 \(V\) 中秩为 \(l\) 的离散子群，并且 \(R\) 的每一个基都是 \(Q(R)\) 的一个基。

同样，\(\mathrm{Q}(\mathsf{R}^{\cdot})\) 是 \(V^{*}\) 中秩为 l 的离散子群。

 命题 26. 对于 \(x \in V\)，把满足 \(\langle x, y^* \rangle \in \mathbf{Z}\) 对所有 \(y^* \in Q(R^*)\)（或等价地，对所有 \(y^* \in R^*)\) 成立的元素所成的集合记为 \(G\)，即 \(V\) 中包含 \(Q(R)\) 的离散子群。若 \(B'\) 是 \(R^*\) 的一个基，则 \(V\) 中与 \(B'\) 对偶的基是 \(G\) 的一个基。

设 \(x \in V\)。以下三个性质等价：

\begin{enumerate}
\def\labelenumi{(\roman{enumi})}
\item
  \(\langle x, y^{*} \rangle \in \mathbf{Z}\) 对所有 \(y^{*} \in \mathrm{Q}(\mathrm{R}^{\vee})\) 成立；
\item
  \(\langle x, y^{*} \rangle \in Z\) 对所有 \(y^{*} \in B'\) 成立；
\item
  \(x\) 关于 \(\mathbf{B}'\) 的对偶基的坐标属于 \(\mathbf{Z}\)。
\end{enumerate}

由此我们得到，与 \(B'\) 对偶的基是 \(G\) 的一个基。另一方面，\((RS_{\mathrm{III}})\) 证明 \(R \subseteq G\)，所以 \(Q(R) \subseteq G\)。

命题 26 中的群 G 记为 P(R)，其元素称为 R 的权。我们也可以考虑 P(R') 这一 R' 的权群。

根据 Algebra, Chap. VII, 2nd edn., § 4, no. 8，

\[
\mathrm{P} (\mathrm{R}) / \mathrm{Q} (\mathrm{R}), \quad \mathrm{P} (\mathrm{R} ^ {\cdot}) / \mathrm{Q} (\mathrm{R} ^ {\cdot})
\]

是关于 \(\mathbf{Q} / \mathbf{Z}\) 对偶的有限群，因而同构。这两个群的公共阶称为 R（或 R \(^{\sim}\)）的连接指标。

若 R 是根系 \(R_{i}\) 的直和，则群 \(Q(R)\)（相应地，\(P(R)\)）典范地等同于 \(Q(R_{i})\)（相应地，\(P(R_{i})\)）的直和。

命题 27. 设 \(\mathbb{R}_1\) 是 \(\mathbb{R}\) 的子集，\(\mathbb{Q}_1\) 是 \(\mathbb{Q}(\mathbb{R})\) 中由 \(\mathbb{R}_1\) 生成的子群，\(\mathbb{W}_1\) 是 \(\mathbb{W}(\mathbb{R})\) 中由 \(s_\alpha\)（\(\alpha \in \mathbb{R}_1\)）生成的子群。若 \(p \in \mathrm{P}(\mathbb{R})\) 且 \(w \in \mathbb{W}_1\)，则 \(p - w(p) \in \mathbb{Q}_1\)。

若 \(w = s_{\alpha}\)，其中 \(\alpha \in \mathbb{R}_1\)，则

\[
p - w (p) = \langle p, \alpha^ {\prime} \rangle \alpha \in \mathbf {Z} \alpha \subseteq \mathrm{Q} _ {1}.
\]

若 \(w = s_{\alpha_1}s_{\alpha_2}\ldots s_{\alpha_r}\)，其中 \(\alpha_{1},\ldots ,\alpha_{r}\in \mathbb{R}_{1}\)，则 \(p - w(p)\in \mathbf{Q}_1\) 仍然成立，正如对 \(r\) 归纳所见。

群 A(R) 保持 P(R) 和 Q(R) 不变，因而作用于商群 P(R)/Q(R)。由命题 27，群 W(R) 在 P(R)/Q(R) 上平凡作用。取商后，我们看到商群 A(R)/W(R)（cf.~命题 16，第 5 号）典范地作用于 P(R)/Q(R)。

